The Mean-Median Difference ($\text{MM} = \bar{v}_D - \tilde{v}_D$), introduced by \citet{wang2016},
measures the gap between the mean and median Democratic vote share across congressional districts.
Positive MM indicates Democratic voters are geographically concentrated in fewer, higher-margin
districts---a pattern associated with packing. Unlike the Efficiency Gap, MM does not require
a swing model and is more stable across election cycles. Its central evidentiary challenge
in litigation is the geographic sorting confound: any positive MM may reflect where Democrats
choose to live, not deliberate mapmaker choices.

We resolve this confound empirically by comparing enacted 2022 congressional maps to
\textsc{bisect} algorithmic maps drawn from the same geographic constraints.
For North Carolina ($k=14$), enacted MM $\approx 0.07$ vs.\ bisect MM $\approx 0.02$
(60--70\% reduction); for Wisconsin ($k=8$), enacted MM $\approx 0.06$ vs.\ bisect MM $\approx 0.02$.
The bisect residual of $\approx 0.02$--$0.03$ establishes the geographic sorting baseline---the
MM that the residential distribution of voters produces in any compact map. The enacted-bisect
gap of $\approx 0.04$--$0.05$ isolates the mapmaker contribution, providing courts with
a precise measure of manipulation independent of residential patterns.

MM is more stable than EG across election cycles (NC standard deviation $\approx 0.015$ vs.\
$\approx 0.030$ for EG), making it the preferred metric for single-election analysis in
post-\textit{Rucho} state court proceedings. It has been cited in NC \textit{Harper v.\ Hall}
(2022) and PA redistricting proceedings; no court has adopted a specific MM threshold.
